\begin{figure}[t]
\begin{center}
\begin{small}
\begin{sc}
\begin{center}
\begin{tikzpicture}
\begin{axis}[
    title={Patch position},
    xlabel={Patch shake (\%)},
    ylabel={Accuracy (\%)},
    xtick={1,3,5,7,9},
    xticklabels={0, 12.5, 25.0, 37.5, 50},
    xmin=1,
    xmax=9,
    ymin=77,
    ymajorgrids=true,
    grid style=dashed,
    scale only axis,
    height=3.8cm,
    width=.8\linewidth,
    legend pos=south west,
    legend cell align={left}
]
    \addplot[color=red]
    coordinates {
    (1,83.800)(2,83.728)(3,83.524)(4,83.222)(5,82.954)(6,82.504)(7,81.964)(8,81.532)(9,80.838)
    };
    \addlegendentry{ViT}

    \addplot[color=darkorange25512714]
    coordinates {
    (1,83.6380022329712)(2,83.6540022796631)(3,83.5340022680664)(4,83.3500022714233)(5,83.0480023071289)(6,82.7980022903442)(7,82.4940022140503)(8,82.1240023007202)(9,81.7680023339844)
    };
    \addlegendentry{MAE}

    \addplot[color=forestgreen4416044]
    coordinates {
    (1,83.7960022872925)(2,82.9260023580933)(3,81.7320023175049)(4,80.1100023526001)(5,78.2520023229981)(6,76.4020022477722)(7,74.1960022163391)(8,71.7180021903992)(9,69.016002142334)
    };
    \addlegendentry{PVT}

    \addplot[color=gray]
    coordinates {
    (1,83.4820023419189)(2,83.2860023699951)(3,83.0220023693848)(4,82.7500023706055)(5,82.2060024127197)(6,81.688002366333)(7,81.0780023422241)(8,80.3400023358154)(9,79.4460023458862)
    };
    \addlegendentry{Swin}
    
    \addplot[color=blue]
    coordinates {
    (1,82.036)(2,81.946)(3,82.012)(4,81.950)(5,81.830)(6,81.738)(7,81.704)(8,81.496)(9,81.294)
    };
    \addlegendentry{\ours{} (ours)}
\end{axis}
\end{tikzpicture}
\end{center}
\end{sc}
\end{small}
\end{center}
\afterplot
\caption{\textbf{Patch position:} The effect of applying positional perturbation on the accuracy. \ours{} can naturally interpret a whole range of possible position values. Being almost immune to the movement of sampled patches it eventually outperforms all but MAE for large perturbations. Note that the X-axis represents the percentage of patch movement relative to the patch size}
\label{fig:shake-sample}
\end{figure}